\documentclass[../../../include/open-logic-section]{subfiles}

\begin{document}

\olfileid{sth}{z}{union}
\olsection{సమ్మేళనం}

\olref[sth][z][sep]{prop:intersectionsexist}
ద్వారా ఛేదనలు లభించాయి.
కానీ ఏవైనా సమితుల
సమ్మేళనాలు ఉండాలని
కోరుకుంటే, మరో స్వీకృతాన్ని
నిర్దేశించాలి:

\begin{axiom}[సమ్మేళనం] ఏ సమితి $A$కైనా,
$\bigcup A =
\Setabs{x}{(\exists b \in A) x \in b}$ సమితి ఉంటుంది.
\[
	\forall A \exists U \forall x(x \in U \liff (\exists b \in A)x \in b)
\]
\end{axiom}

సంచిత-దశలవారీ భావన
ఈ స్వీకృతాన్ని కూడా
సమర్థిస్తుంది. $A$ ఒక
సమితి అనుకుందాం; అప్పుడు
\stageshier\ ప్రకారం $A$
ఏదో ఒక దశ $S$లో
ఏర్పడుతుంది. $A$లోని
ప్రతి మూలకమూ $S$కు
\emph{ముందు} ఏర్పడింది
(\stagesacc\ ప్రకారం);
అందువల్ల అదే తర్కంతో,
$A$లోని ప్రతి మూలకపు
ప్రతి మూలకమూ $S$కు
ముందే ఏర్పడింది. కాబట్టి
ఆ సమితులన్నీ $S$కు
ముందే అందుబాటులో ఉన్నాయి;
వాటిని $S$లో ఒక సమితిగా
ఏర్పరచవచ్చు. ఆ సమితియే
$\bigcup A$.

\end{document}
